\section{Contrastive Probe Validation}

\subsection{Purpose}

This section moves from a probe queue to controlled decipherment constraints. The previous model listed promising minimal-pair neighborhoods. This validation layer asks which probes can already be promoted to grammar constraints, which can become phonetic probes after image checking, and which must remain controls.

The key rule remains strict: no sound value, lexical translation, or language-family claim is allowed unless the proposed value predicts all occurrences of the relevant contrast.

\subsection{Method}

The script \texttt{scripts/contrastive-probe-validation.py} combines:

\begin{itemize}
  \item the top thirty minimal-pair probes;
  \item exact occurrence profiles from structural reconstructions;
  \item same-frame minimal-pair counts;
  \item component roles from the morpheme assignment table;
  \item site, object, material, symbol, and semantic-frame distributions.
\end{itemize}

The validator separates three levels of result:

\begin{itemize}
  \item grammar constraints, which can be used now;
  \item stem-semantic probes, which are promising but occurrence-limited;
  \item phonetic-ready probes, which still require image checking before any abstract sound contrast is proposed.
\end{itemize}

\subsection{Validation Result}

\begin{table}[htbp]
\centering
\begin{tabular}{lr}
\toprule
\textbf{Metric} & \textbf{Value} \\
\midrule
Probe validations & 30 \\
Morphology constraints & 11 \\
Terminal constraints & 1 \\
Phonetic-ready after image check & 1 \\
Stem-semantic probes & 3 \\
Full-neighborhood role-boundary controls & 108 \\
Stem lattice edges & 7 \\
\bottomrule
\end{tabular}
\caption{Contrastive probe validation summary.}
\label{tab:contrastive-probe-validation-summary}
\end{table}

\subsection{First Phonetic-Ready Edge}

The first probe to reach phonetic-ready status after image checking is:

\begin{center}
\texttt{240-482 :: 240-904}
\end{center}

in the frame:

\begin{center}
\texttt{002|740}
\end{center}

The controlled interpretation is:

\begin{center}
\texttt{240 + STEM}
\end{center}

where \texttt{482} and \texttt{904} occupy the same second-position stem slot. The current evidence is still pre-phonetic: it licenses a future contrast test, not a sound value.

\subsection{Stem Lattice}

The strongest stem-lattice edges are:

\begin{table}[htbp]
\centering
\footnotesize
\begin{tabular}{lp{0.36\textwidth}lr}
\toprule
\textbf{Edge} & \textbf{Probe} & \textbf{Status} & \textbf{Score} \\
\midrule
\texttt{176:048} & \texttt{240-176 :: 048} & Stem-semantic probe & 0.822 \\
\texttt{482:904} & \texttt{240-482 :: 240-904} & Phonetic-ready after image check & 0.819 \\
\texttt{176:048} & \texttt{176 :: 048} & Stem-semantic probe & 0.788 \\
\texttt{176:061} & \texttt{176 :: 061} & Stem-semantic probe & 0.750 \\
\texttt{176:482} & \texttt{240-176 :: 240-482} & Hold for more occurrences & 0.698 \\
\texttt{482:773} & \texttt{240-482 :: 240-773} & Hold for more occurrences & 0.696 \\
\bottomrule
\end{tabular}
\caption{Top stem-lattice contrasts.}
\label{tab:stem-lattice-contrasts}
\end{table}

\subsection{Grammar Constraints Now Licensed}

The validation layer licenses the following grammar constraints:

\begin{itemize}
  \item \texttt{240} behaves as a qualifier/title/classifier-like component in several neighborhoods.
  \item \texttt{240} and \texttt{235} compete in title/classifier/formula-opener environments.
  \item \texttt{740} and \texttt{520} form a terminal/suffix contrast family.
  \item \texttt{032} must remain blocked for phonetics because it creates many role-boundary controls.
\end{itemize}

\subsection{Controlled Claim Ledger}

The current claim ledger now contains six controlled claims:

\begin{itemize}
  \item \texttt{C-240-QUALIFIER}: supported morphology constraint.
  \item \texttt{C-176-STEM}: stem-semantic probe.
  \item \texttt{C-482-STEM-SLOT}: emerging stem lattice.
  \item \texttt{C-740-520-TERMINAL}: terminal constraint.
  \item \texttt{C-032-CONTROL}: blocked for phonetics.
  \item \texttt{C-LANGUAGE-GATE}: language identification remains blocked.
\end{itemize}

\subsection{Interpretation}

This is a meaningful step closer to decipherment because the project now has its first phonetic-ready contrast candidate:

\begin{center}
\texttt{482 :: 904}
\end{center}

inside:

\begin{center}
\texttt{240 + STEM}
\end{center}

The next safe step is visual validation of every \texttt{240-482} and \texttt{240-904} occurrence. If the visual and contextual evidence holds, the project can begin testing an abstract contrast between two stem values. That would still not be a language reading, but it would be the first responsible opening toward phonetic decipherment.

\subsection{Outputs}

\begin{itemize}
  \item \texttt{outputs/contrastive\_probe\_validation\_summary.csv}
  \item \texttt{outputs/validated\_probe\_results.csv}
  \item \texttt{outputs/probe\_occurrence\_contexts.csv}
  \item \texttt{outputs/stem\_contrast\_lattice.csv}
  \item \texttt{outputs/controlled\_decipherment\_claims.csv}
  \item \texttt{outputs/contrastive\_probe\_validation.tex}
\end{itemize}
